\scanpage{209}
\mspage{83}
limites inductives filtrantes, $C^{\unsure{\pi|R}}=C\cap\mathrm{Filt}^\pi$ est stable par $\underrightarrow{\lim}$ filtrantes de cardinal $\leq\pi$, en vertu de la condition d) sur $\mathrm{Filt}^\pi$.

Considérons maintenant le carré cocartésien
\diagram{\begin{tikzcd}[ampersand replacement=\&]
 A'\arrow[r,hook',"i'"]\arrow[d,hook]\&B'\arrow[d]\&\\
 A\arrow[r,hook',"{\deleted{\unsure{g}}i}"]\& C'=A\amalg_{A'}B'\arrow[r,"{\unsure{j|J}}"]\&B
\end{tikzcd}}
Comme $i'\in C^\pi$, on a $i\in\widetilde{C^\pi}$ (car $\widetilde{C^\pi}$ stable par cochgt de base). D'autre part, j'ai besoin du fait que $j$ soit dans $C$, \deleted{\unsure{a priori} \unsure{plu}} et je \uline{ne} sais pas $=$ qu'il est mono ! \unsure{Aussi} il faut repasser aux notations indicielles
\diagram{\begin{tikzcd}[ampersand replacement=\&]
 A_i\arrow[r,hook,no head,"\alpha_i"]\arrow[d,hook,no head]\&B_i\arrow[d]\&\\
 A\arrow[r,hook,"\beta_i"]\&C_i\arrow[r,"{\unsure{\gamma_i|\delta_i}}"]\&B
\end{tikzcd}}
Comme \unsure{le foncteur|la formation} des sommes amalgamées commute aux limites inductives filtrantes, on sait que
\[
 \begin{aligned}
 \unsure{\underline{\lim}|\underrightarrow{\lim}}\,^{\unsure{\prime\prime}} C_i
 &=\deleted{\underrightarrow{\lim}A_i}\ A\amalg_{\underrightarrow{\lim} A_i}\underline{\lim}B_i\\
 &=A\amalg_A B=B,
 \end{aligned}
\]
donc la limite inductive des $\gamma_i$ est un iso. (\unsure{can.} $\mathrm{id}_B$). \unsure{D'autre part}, le $=$ argument montre que le système inductif des $C_i$ $(i\in J)$ \inserted{de $M$} est \deleted{lui aussi} $\pi$-adapté, donc aussi celui des $\gamma_i$.$^{\ast}$

% slanted
\margin{\uuline{Sauf} qu'on ne sait pas que les $C_i$ soient dans $\mathrm{Filt}^\pi$ --- en fait ils ne seront \unsure{guère|pas} en général, \illegible{2} \unsure{seraient|serait} \unsure{supposer|supposés}, \unsure{puisque} $A,B\in\mathrm{Filt}^{\unsure{\pi}}$}
% archive sheet number: 208
